\documentclass[12pt, a4paper]{article}

% ── Pacotes ──────────────────────────────────────────────────────────────────
\usepackage[utf8]{inputenc}
\usepackage[T1]{fontenc}
\usepackage[brazil]{babel}
\usepackage{geometry}
\usepackage{graphicx}
\usepackage{hyperref}
\usepackage{listings}
\usepackage{xcolor}
\usepackage{booktabs}
\usepackage{array}
\usepackage{tabularx}
\usepackage{amsmath}
\usepackage{fancyhdr}
\usepackage{titlesec}
\usepackage{setspace}
\usepackage{parskip}
\usepackage{tcolorbox}
\usepackage{enumitem}
\usepackage{microtype}
\usepackage{lmodern}
\usepackage{caption}
\usepackage{subcaption}
\usepackage{float}

% ── Geometria ────────────────────────────────────────────────────────────────
\geometry{
  top=3cm, bottom=2cm,
  left=3cm, right=2cm,
  headheight=15pt
}

% ── Cores ────────────────────────────────────────────────────────────────────
\definecolor{accent}{HTML}{00875A}
\definecolor{codebg}{HTML}{F4F4F4}
\definecolor{codeframe}{HTML}{CCCCCC}
\definecolor{keyword}{HTML}{0033B3}
\definecolor{string}{HTML}{067D17}
\definecolor{comment}{HTML}{8C8C8C}
\definecolor{number}{HTML}{1750EB}

% ── Hyperref ─────────────────────────────────────────────────────────────────
\hypersetup{
  colorlinks=true,
  linkcolor=accent,
  urlcolor=accent,
  citecolor=accent,
  pdftitle={Calculadora API com FastAPI},
  pdfauthor={Gabriel Alvarenga; Rafael Fernandes}
}

% ── Listings (código) ────────────────────────────────────────────────────────
\lstdefinestyle{python}{
  language=Python,
  backgroundcolor=\color{codebg},
  frame=single,
  rulecolor=\color{codeframe},
  framerule=0.5pt,
  basicstyle=\ttfamily\small,
  keywordstyle=\color{keyword}\bfseries,
  stringstyle=\color{string},
  commentstyle=\color{comment}\itshape,
  numberstyle=\tiny\color{comment},
  numbers=left,
  numbersep=8pt,
  stepnumber=1,
  showstringspaces=false,
  breaklines=true,
  breakatwhitespace=true,
  tabsize=4,
  captionpos=b,
  literate=
    {á}{{\'a}}1 {é}{{\'e}}1 {í}{{\'i}}1 {ó}{{\'o}}1 {ú}{{\'u}}1
    {Á}{{\'A}}1 {É}{{\'E}}1 {Í}{{\'I}}1 {Ó}{{\'O}}1 {Ú}{{\'U}}1
    {ã}{{\~a}}1 {õ}{{\~o}}1 {Ã}{{\~A}}1 {Õ}{{\~O}}1
    {â}{{\^a}}1 {ê}{{\^e}}1 {ô}{{\^o}}1
    {ç}{{\c{c}}}1 {Ç}{{\c{C}}}1
}

\lstdefinestyle{bash}{
  language=bash,
  backgroundcolor=\color{black!90},
  frame=single,
  rulecolor=\color{black},
  framerule=0.5pt,
  basicstyle=\ttfamily\small\color{white},
  keywordstyle=\color{green!60!white}\bfseries,
  stringstyle=\color{yellow!80!white},
  commentstyle=\color{gray!60},
  showstringspaces=false,
  breaklines=true,
  tabsize=2,
  captionpos=b
}

\lstset{style=python}

% ── Cabeçalho/rodapé ─────────────────────────────────────────────────────────
\pagestyle{fancy}
\fancyhf{}
\fancyhead[L]{\small\textcolor{accent}{\textbf{Calculadora API}} — Sistemas Distribuídos}
\fancyhead[R]{\small UNILAGO}
\fancyfoot[C]{\small\thepage}
\renewcommand{\headrulewidth}{0.5pt}

% ── Caixas coloridas ─────────────────────────────────────────────────────────
\tcbuselibrary{skins, breakable}
\newtcolorbox{infobox}[1][]{
  colback=accent!8,
  colframe=accent,
  fonttitle=\bfseries,
  title=#1,
  breakable
}

% ── Títulos ──────────────────────────────────────────────────────────────────
\titleformat{\section}{\Large\bfseries\color{accent}}{}{0em}{\thesection\quad}
\titleformat{\subsection}{\large\bfseries}{}{0em}{\thesubsection\quad}
\titleformat{\subsubsection}{\normalsize\bfseries}{}{0em}{\thesubsubsection\quad}

% ─────────────────────────────────────────────────────────────────────────────
\begin{document}

% ── Capa ─────────────────────────────────────────────────────────────────────
\begin{titlepage}
  \centering
  \vspace*{2cm}

  {\large Universidade do Grandes Lagos — UNILAGO}\\[0.3cm]
  {\large Curso de Ciência da Computação}\\[0.3cm]
  {\large Disciplina: Programação de Sistemas Distribuídos}

  \vspace{2cm}
  \rule{\linewidth}{2pt}\\[0.4cm]
  {\Huge\bfseries\color{accent} Calculadora API}\\[0.3cm]
  {\Large com FastAPI e SQLite}\\[0.2cm]
  \rule{\linewidth}{2pt}

  \vspace{1.5cm}

  \begin{tabular}{rl}
    \textbf{Discentes:} & Gabriel Henrique Mendes Angelo Alvarenga \\
                        & Rafael Martins Fernandes \\[0.6cm]
    \textbf{Professor:} & Eng. Gleydes Oliveira \\[0.6cm]
    \textbf{Versão da API:} & 2.0.0 \\[0.3cm]
    \textbf{Data:}      & \today
  \end{tabular}

  \vfill
  {\small São José do Rio Preto — SP}
\end{titlepage}

% ── Resumo ───────────────────────────────────────────────────────────────────
\begin{abstract}
\noindent
Este artigo documenta o desenvolvimento de uma API RESTful de calculadora utilizando o framework FastAPI em Python, como projeto avaliativo da disciplina de Programação de Sistemas Distribuídos. A solução implementa as quatro operações matemáticas básicas — soma, subtração, multiplicação e divisão —, além de potenciação e raiz quadrada. Como evolução em relação ao nível básico proposto, foram incorporados recursos de nível intermediário: persistência de histórico de cálculos por usuário com banco de dados SQLite, habilitação de CORS para comunicação com frontend, e um total de seis operações disponíveis. A interface web foi desenvolvida em HTML e JavaScript puro, consumindo três ou mais endpoints da API. A documentação automática via Swagger UI e o tratamento adequado de códigos de status HTTP são também abordados.

\vspace{0.5em}
\noindent\textbf{Palavras-chave:} FastAPI, API REST, Sistemas Distribuídos, SQLite, Python, CORS.
\end{abstract}

\newpage
\tableofcontents
\newpage

% ─────────────────────────────────────────────────────────────────────────────
\section{Introdução}

Em sistemas distribuídos, a comunicação entre componentes independentes é um dos desafios centrais da engenharia de software. APIs (\textit{Application Programming Interfaces}) do estilo REST (\textit{Representational State Transfer}) tornaram-se o padrão de facto para essa comunicação, oferecendo simplicidade, independência de linguagem e escalabilidade.

Este projeto tem por objetivo a construção de uma API de calculadora, partindo do nível básico especificado pelo professor e avançando para o nível intermediário com a adição de:

\begin{itemize}[leftmargin=*]
  \item Histórico de cálculos persistido por usuário em banco de dados SQLite;
  \item CORS habilitado para consumo pelo frontend;
  \item Duas operações adicionais: potenciação (\texttt{POST /potencia}) e raiz quadrada (\texttt{POST /raiz});
  \item Frontend em HTML + JavaScript puro que consome a API.
\end{itemize}

O desenvolvimento foi realizado com Python 3.11, FastAPI 0.115, Uvicorn 0.30 e Pydantic 2.9.

% ─────────────────────────────────────────────────────────────────────────────
\section{Fundamentos Teóricos}

\subsection{API REST}

REST é um estilo arquitetural descrito por Roy Fielding em sua dissertação de 2000. Uma API RESTful utiliza os verbos HTTP de forma semântica:

\begin{table}[H]
  \centering
  \caption{Verbos HTTP e seus usos}
  \begin{tabularx}{\linewidth}{lX}
    \toprule
    \textbf{Verbo} & \textbf{Uso} \\
    \midrule
    GET    & Consultar dados sem modificar o servidor \\
    POST   & Enviar dados para processamento \\
    PUT    & Atualizar um recurso completo \\
    DELETE & Remover um recurso \\
    \bottomrule
  \end{tabularx}
\end{table}

\subsection{FastAPI}

FastAPI é um framework Python moderno e de alto desempenho para construção de APIs. Seus principais diferenciais são: geração automática de documentação Swagger e ReDoc; validação de dados via Pydantic; suporte nativo a programação assíncrona (async/await); e desempenho comparável a servidores escritos em Go e Node.js.

\subsection{Pydantic}

Pydantic é a biblioteca de validação de dados utilizada pelo FastAPI. Com ela, define-se modelos (classes Python com tipos anotados) que descrevem a estrutura dos dados. Quando o cliente envia dados inválidos, o FastAPI retorna automaticamente o código HTTP 422 com detalhes precisos do erro — sem necessidade de código de validação manual.

\subsection{SQLite}

SQLite é um banco de dados relacional embutido, baseado em um único arquivo. É ideal para projetos de pequeno e médio porte que não necessitam de um servidor de banco de dados dedicado. Neste projeto, o SQLite é utilizado para persistir o histórico de cálculos por usuário.

\subsection{CORS}

CORS (\textit{Cross-Origin Resource Sharing}) é um mecanismo de segurança dos navegadores que impede requisições a domínios diferentes do da página origem. Para que o frontend (servido localmente) possa consumir a API, é necessário configurar os cabeçalhos CORS no servidor.

% ─────────────────────────────────────────────────────────────────────────────
\section{Configuração do Ambiente}

\subsection{Pré-requisitos}

\begin{itemize}[leftmargin=*]
  \item Python 3.8 ou superior;
  \item pip (gerenciador de pacotes Python).
\end{itemize}

\subsection{Criação do Ambiente Virtual}

\begin{lstlisting}[style=bash, caption={Criação e ativação do ambiente virtual}]
# Criar pasta do projeto
mkdir calculadora-api && cd calculadora-api

# Criar ambiente virtual
python -m venv venv

# Ativar (Linux/macOS)
source venv/bin/activate

# Ativar (Windows)
venv\Scripts\activate
\end{lstlisting}

\subsection{Instalação das Dependências}

O arquivo \texttt{requirements.txt} declara as versões fixas de cada pacote:

\begin{lstlisting}[style=bash]
fastapi==0.115.0
uvicorn==0.30.6
pydantic==2.9.2
\end{lstlisting}

\begin{lstlisting}[style=bash]
pip install -r requirements.txt
\end{lstlisting}

% ─────────────────────────────────────────────────────────────────────────────
\section{Estrutura do Projeto}

\begin{lstlisting}[style=bash, caption={Estrutura de arquivos do projeto}]
calculadora-api/
|-- venv/              # Ambiente virtual (NAO commitar!)
|-- main.py            # API principal
|-- requirements.txt   # Dependencias
|-- test_client.py     # Script de testes Python
|-- calculadora.db     # Banco SQLite (gerado automaticamente)
|-- frontend/
|   +-- index.html     # Interface web
+-- README.md
\end{lstlisting}

% ─────────────────────────────────────────────────────────────────────────────
\section{Implementação da API}

\subsection{Inicialização e CORS}

A instância do FastAPI é criada com metadados que aparecem automaticamente na documentação Swagger. O middleware de CORS é adicionado logo após, permitindo que qualquer origem consuma a API durante o desenvolvimento:

\begin{lstlisting}[caption={Criação do app e configuração do CORS}]
from fastapi import FastAPI, HTTPException, Depends
from fastapi.middleware.cors import CORSMiddleware

app = FastAPI(
    title="Calculadora API",
    description="API para Sistemas Distribuidos",
    version="2.0.0",
)

app.add_middleware(
    CORSMiddleware,
    allow_origins=["*"],
    allow_credentials=True,
    allow_methods=["*"],
    allow_headers=["*"],
)
\end{lstlisting}

\subsection{Banco de Dados SQLite}

O banco é criado automaticamente na primeira execução. A função \texttt{get\_db()} é utilizada como dependência nos endpoints via \texttt{Depends()}:

\begin{lstlisting}[caption={Inicialização do SQLite e injeção de dependência}]
import sqlite3

DB_PATH = "calculadora.db"

def get_db():
    conn = sqlite3.connect(DB_PATH)
    conn.row_factory = sqlite3.Row
    try:
        yield conn
    finally:
        conn.close()

def init_db():
    conn = sqlite3.connect(DB_PATH)
    conn.execute("""
        CREATE TABLE IF NOT EXISTS historico (
            id        INTEGER PRIMARY KEY AUTOINCREMENT,
            usuario   TEXT    NOT NULL,
            operacao  TEXT    NOT NULL,
            numero1   REAL    NOT NULL,
            numero2   REAL,
            resultado REAL    NOT NULL,
            data_hora TEXT    NOT NULL
        )
    """)
    conn.commit()
    conn.close()

init_db()
\end{lstlisting}

\subsection{Modelos Pydantic}

São definidos três modelos: \texttt{OperacaoRequest} (entrada binária), \texttt{OperacaoUnariaRequest} (entrada unária, para raiz) e \texttt{ResultadoResponse} (saída):

\begin{lstlisting}[caption={Modelos Pydantic de entrada e saída}]
from pydantic import BaseModel
from typing import Optional

class OperacaoRequest(BaseModel):
    numero1: float
    numero2: float
    usuario: str = "anonimo"

class OperacaoUnariaRequest(BaseModel):
    numero1: float
    usuario: str = "anonimo"

class ResultadoResponse(BaseModel):
    operacao: str
    numero1: float
    numero2: Optional[float] = None
    resultado: float
    usuario: str
    data_hora: str
\end{lstlisting}

\subsection{Endpoints de Cálculo}

Cada endpoint segue o mesmo padrão: recebe os dados validados pelo Pydantic, realiza o cálculo, persiste no banco via \texttt{salvar\_historico()} e retorna um \texttt{ResultadoResponse}. O exemplo abaixo mostra a divisão, que inclui validação extra:

\begin{lstlisting}[caption={Endpoint de divisão com tratamento de erro}]
@app.post("/dividir", response_model=ResultadoResponse)
def dividir(dados: OperacaoRequest, conn=Depends(get_db)):
    if dados.numero2 == 0:
        raise HTTPException(
            status_code=400,
            detail="Divisao por zero nao e permitida!"
        )
    resultado = dados.numero1 / dados.numero2
    data_hora = salvar_historico(
        conn, dados.usuario, "divisao",
        dados.numero1, resultado, dados.numero2
    )
    return ResultadoResponse(
        operacao="divisao",
        numero1=dados.numero1,
        numero2=dados.numero2,
        resultado=resultado,
        usuario=dados.usuario,
        data_hora=data_hora,
    )
\end{lstlisting}

\subsection{Raiz Quadrada}

A raiz quadrada é uma operação unária (recebe apenas um número) e valida se o número é não-negativo:

\begin{lstlisting}[caption={Endpoint de raiz quadrada}]
import math

@app.post("/raiz", response_model=ResultadoResponse)
def raiz_quadrada(dados: OperacaoUnariaRequest, conn=Depends(get_db)):
    if dados.numero1 < 0:
        raise HTTPException(
            status_code=400,
            detail="Raiz de numero negativo nao e permitida!"
        )
    resultado = math.sqrt(dados.numero1)
    data_hora = salvar_historico(
        conn, dados.usuario, "raiz_quadrada",
        dados.numero1, resultado
    )
    return ResultadoResponse(
        operacao="raiz_quadrada",
        numero1=dados.numero1,
        resultado=resultado,
        usuario=dados.usuario,
        data_hora=data_hora,
    )
\end{lstlisting}

\subsection{Histórico por Usuário}

O endpoint abaixo retorna todos os cálculos de um determinado usuário, ordenados do mais recente ao mais antigo:

\begin{lstlisting}[caption={Endpoint de histórico por usuário}]
@app.get("/historico/{usuario}")
def historico_usuario(usuario: str, conn=Depends(get_db)):
    rows = conn.execute(
        "SELECT * FROM historico WHERE usuario = ? ORDER BY id DESC",
        (usuario,),
    ).fetchall()
    if not rows:
        raise HTTPException(
            status_code=404,
            detail=f"Nenhum historico encontrado para '{usuario}'."
        )
    return {
        "usuario": usuario,
        "total": len(rows),
        "historico": [dict(r) for r in rows]
    }
\end{lstlisting}

% ─────────────────────────────────────────────────────────────────────────────
\section{Resumo dos Endpoints}

\begin{table}[H]
  \centering
  \caption{Todos os endpoints da API v2.0}
  \begin{tabularx}{\linewidth}{llX}
    \toprule
    \textbf{Método} & \textbf{Rota} & \textbf{Descrição} \\
    \midrule
    GET  & \texttt{/}                    & Boas-vindas \\
    POST & \texttt{/somar}               & Soma dois números \\
    POST & \texttt{/subtrair}            & Subtrai dois números \\
    POST & \texttt{/multiplicar}         & Multiplica dois números \\
    POST & \texttt{/dividir}             & Divide (valida divisão por zero) \\
    POST & \texttt{/potencia}            & Calcula \texttt{numero1\^{}numero2} \\
    POST & \texttt{/raiz}                & Raiz quadrada de \texttt{numero1} \\
    GET  & \texttt{/calcular}            & Operação via \textit{query parameters} \\
    GET  & \texttt{/historico/\{usuario\}} & Histórico filtrado por usuário \\
    GET  & \texttt{/historico}           & Todo o histórico \\
    \bottomrule
  \end{tabularx}
\end{table}

% ─────────────────────────────────────────────────────────────────────────────
\section{Códigos de Status HTTP}

\begin{table}[H]
  \centering
  \caption{Códigos de status retornados pela API}
  \begin{tabularx}{\linewidth}{llX}
    \toprule
    \textbf{Código} & \textbf{Nome} & \textbf{Quando ocorre} \\
    \midrule
    200 & OK                    & Operação realizada com sucesso \\
    400 & Bad Request           & Divisão por zero, raiz de negativo ou operação inválida \\
    404 & Not Found             & Usuário sem histórico ou endpoint inexistente \\
    422 & Unprocessable Entity  & Dado inválido (ex.: texto no lugar de número) \\
    500 & Internal Server Error & Erro inesperado no servidor \\
    \bottomrule
  \end{tabularx}
\end{table}

% ─────────────────────────────────────────────────────────────────────────────
\section{Frontend}

A interface web foi desenvolvida em HTML5 e JavaScript puro (sem frameworks), consumindo os seguintes endpoints da API:

\begin{enumerate}[leftmargin=*]
  \item \texttt{POST /somar}, \texttt{/subtrair}, \texttt{/multiplicar}, \texttt{/dividir}, \texttt{/potencia}, \texttt{/raiz} — para realizar os cálculos;
  \item \texttt{GET /historico/\{usuario\}} — para exibir o histórico filtrado por usuário;
  \item \texttt{GET /historico} — para exibir o histórico geral.
\end{enumerate}

A comunicação é feita via \texttt{fetch()} com método \texttt{POST} e \texttt{Content-Type: application/json}. Os resultados são exibidos dinamicamente na página sem recarregamento.

O frontend apresenta:
\begin{itemize}[leftmargin=*]
  \item Campo de identificação de usuário;
  \item Seleção visual da operação por \textit{chips} clicáveis;
  \item Campos numéricos com ocultamento automático do segundo número para operações unárias;
  \item Exibição do resultado com metadados (expressão, usuário, data/hora);
  \item Painel lateral de histórico com busca por usuário e listagem de todos os registros.
\end{itemize}

% ─────────────────────────────────────────────────────────────────────────────
\section{Testes}

\subsection{Swagger UI}

Após iniciar o servidor com \texttt{uvicorn main:app --reload}, o Swagger UI é acessível em \url{http://localhost:8000/docs}. A interface permite testar todos os endpoints interativamente no navegador.

\subsection{cURL}

Exemplos de teste via terminal:

\begin{lstlisting}[style=bash, caption={Exemplos de testes com cURL}]
# Soma
curl -X POST "http://localhost:8000/somar" \
     -H "Content-Type: application/json" \
     -d '{"numero1": 10, "numero2": 5, "usuario": "gabriel"}'

# Divisao por zero (erro 400 esperado)
curl -X POST "http://localhost:8000/dividir" \
     -H "Content-Type: application/json" \
     -d '{"numero1": 10, "numero2": 0, "usuario": "gabriel"}'

# Historico do usuario
curl "http://localhost:8000/historico/gabriel"
\end{lstlisting}

\subsection{Script Python}

O arquivo \texttt{test\_client.py} automatiza os testes básicos usando apenas a biblioteca padrão (\texttt{urllib}), sem dependências externas, e imprime os resultados no terminal.

% ─────────────────────────────────────────────────────────────────────────────
\section{Contexto em Sistemas Distribuídos}

\subsection{Comunicação Síncrona}

A calculadora utiliza comunicação síncrona: o cliente envia a requisição e aguarda a resposta antes de continuar. Esse modelo é adequado para operações rápidas como cálculos matemáticos. Em sistemas distribuídos de maior escala, operações demoradas são tratadas com comunicação assíncrona via filas de mensagens (ex.: RabbitMQ, Apache Kafka).

\subsection{Baixo Acoplamento}

A separação entre frontend e backend via API REST é um exemplo prático de baixo acoplamento: o frontend não conhece os detalhes internos da API, e a API não conhece os detalhes do frontend. Qualquer cliente que entenda HTTP pode consumir a API — seja um navegador, um \textit{script} Python, o Insomnia ou um serviço escrito em Java.

\subsection{Servidor ASGI}

O Uvicorn atua como servidor ASGI (\textit{Asynchronous Server Gateway Interface}), fazendo a ponte entre o sistema operacional e o FastAPI. Em ambiente de produção, o Uvicorn seria colocado atrás de um proxy reverso como o Nginx para lidar com múltiplas conexões simultâneas e fornecer TLS.

\subsection{Escalabilidade}

Uma API RESTful como a desenvolvida pode ser escalada horizontalmente: múltiplas instâncias do servidor podem ser executadas em paralelo atrás de um balanceador de carga (\textit{load balancer}), com o banco de dados compartilhado entre elas. Em um cenário de microsserviços, a calculadora seria apenas um dos vários serviços independentes que compõem o sistema.

% ─────────────────────────────────────────────────────────────────────────────
\section{Evolução: Básico → Intermediário}

A tabela abaixo compara os recursos implementados no nível básico (conforme especificado no projeto) e os recursos adicionados no nível intermediário:

\begin{table}[H]
  \centering
  \caption{Comparação entre os níveis básico e intermediário}
  \begin{tabularx}{\linewidth}{lXX}
    \toprule
    \textbf{Recurso} & \textbf{Nível Básico} & \textbf{Nível Intermediário} \\
    \midrule
    Banco de dados    & Não             & SQLite (arquivo local) \\
    Histórico         & Não             & Sim, por usuário autenticado \\
    Resposta inclui   & Apenas resultado & Resultado + usuário + data/hora \\
    Operações         & 4 (básicas)      & 6 (+ potência e raiz) \\
    CORS              & Não             & Habilitado para frontend \\
    Frontend          & Não             & HTML + JS puro \\
    \bottomrule
  \end{tabularx}
\end{table}

% ─────────────────────────────────────────────────────────────────────────────
\section{Conclusão}

O projeto demonstrou na prática os conceitos fundamentais de APIs RESTful em um contexto de sistemas distribuídos. A utilização do FastAPI permitiu criar uma API robusta com validação automática de dados, documentação interativa e alto desempenho, com pouco código.

A adição de recursos do nível intermediário — histórico persistido em SQLite, CORS e operações extras — tornou a solução mais realista e representativa dos sistemas distribuídos modernos. O frontend em HTML e JavaScript puro evidenciou o baixo acoplamento característico da arquitetura REST: o cliente e o servidor são desenvolvidos e implantados de forma completamente independente.

Como trabalhos futuros, sugere-se a adição de autenticação JWT, migração do SQLite para um banco de dados relacional como PostgreSQL, e a implantação da API em uma plataforma de nuvem (ex.: Render, Railway ou AWS).

% ─────────────────────────────────────────────────────────────────────────────
\section*{Referências}

\begin{enumerate}[leftmargin=*, label={[\arabic*]}]
  \item FASTAPI. \textit{FastAPI Documentation}. Disponível em: \url{https://fastapi.tiangolo.com}. Acesso em: mar. 2026.
  \item PYDANTIC. \textit{Pydantic v2 Documentation}. Disponível em: \url{https://docs.pydantic.dev}. Acesso em: mar. 2026.
  \item UVICORN. \textit{Uvicorn — ASGI server}. Disponível em: \url{https://www.uvicorn.org}. Acesso em: mar. 2026.
  \item FIELDING, R. T. \textit{Architectural Styles and the Design of Network-based Software Architectures}. Dissertação (Doutorado) — University of California, Irvine, 2000.
  \item SQLITE. \textit{SQLite Documentation}. Disponível em: \url{https://www.sqlite.org/docs.html}. Acesso em: mar. 2026.
  \item MDN WEB DOCS. \textit{Cross-Origin Resource Sharing (CORS)}. Disponível em: \url{https://developer.mozilla.org/en-US/docs/Web/HTTP/CORS}. Acesso em: mar. 2026.
\end{enumerate}

\end{document}
